\section{Visuomotor Congruence and Synchrony}
\label{sec:background_visuomotor_congruence}

Visuomotor congruence describes the correspondence between a user's physical movement and the movement they see. In virtual embodiment research, this correspondence is one of the central cues linking the physical body to the virtual body. When the user performs an action and the virtual body moves in a spatially and temporally corresponding way, the virtual movement can be interpreted as the visual consequence of the user's own action.
% TODO cite: visuomotor congruence / virtual embodiment source.
% TODO cite: agency and body ownership under active movement.

Two aspects of congruence are especially relevant: temporal synchrony and spatial correspondence. Temporal synchrony concerns whether the virtual movement occurs at the same time as the user's real movement. Even small delays can affect the perceived connection between action and feedback, particularly for agency. Spatial correspondence concerns whether the virtual movement occurs in the expected position, direction, and configuration relative to the user's body. A movement can be temporally synchronous but spatially transformed, or spatially aligned but delayed.
% TODO cite: temporal delay / agency threshold studies.
% TODO cite: spatial congruence / body ownership / virtual hand literature.

The distinction between temporal and spatial congruence matters for VR rowing. A boat-centric representation may respond immediately to the user's stroke, preserving temporal synchrony, while still transforming the spatial path of the hands into an oar-based movement. An ergometer-congruent representation may preserve both the timing and the approximate spatial logic of the handle pull, because the virtual hands and handle can follow a path closer to the real ergometer movement. Both representations may therefore be responsive, but they differ in how much of the user's physical movement geometry is preserved.

Classic body illusion research often emphasizes the importance of synchronous multisensory stimulation. In the rubber hand illusion, synchronous visual and tactile cues can support ownership over an artificial hand, while asynchronous cues tend to weaken the illusion. In active virtual hand paradigms, the relationship between the user's own movement and the seen movement becomes especially important, because movement itself provides a continuous cue linking the physical and virtual body.
% TODO cite: rubber hand illusion synchrony.
% TODO cite: moving rubber hand / virtual hand illusion / active movement studies.
% TODO cite: Kalckert and Ehrsson or similar moving rubber hand work if used.

In VR rowing, the user's hand movement is not isolated from the rest of the body. The stroke is a coordinated sequence involving leg drive, trunk swing, arm pull, handle movement, and seat motion. This makes congruence more complex than simply matching a virtual hand to a real hand. A system can align some cues while transforming others. For example, the timing of the stroke may match the user, while the visible trajectory of the virtual hands may follow an oar-based path rather than the real ergometer handle path. Conversely, an ergometer-congruent representation can align the hand and handle more closely while sacrificing some of the visual realism associated with rowing on water.

This is the core reason why the comparison in this thesis is not a simple test of realism. A realistic-looking boat scene may provide strong contextual plausibility, but still introduce a visuomotor transformation between the felt and seen action. A visually less sport-realistic ergometer representation may provide stronger congruence with the user's physical action. The two representation strategies therefore emphasize different forms of fidelity.

The expected role of visuomotor congruence is especially strong for sense of agency. If the virtual action closely follows the user's intended and performed movement, users may experience the action as more directly their own. Body ownership may also benefit from congruence, especially when the virtual hands and body appear spatially aligned with the felt body. Self-location may be affected more indirectly, depending on whether the user experiences themselves as situated in the virtual body, the physical ergometer setup, or the represented rowing environment.
% TODO cite: agency and congruent action feedback.
% TODO cite: ownership and active movement congruence.
% TODO cite: self-location and multisensory body alignment.

At the same time, congruence should not be treated as the only value in the design of VR sport systems. A perfectly congruent representation of the ergometer may not be the representation users find most meaningful as rowing. A transformed boat-centric representation may reduce spatial congruence while increasing perceived sport realism. This thesis therefore examines visuomotor congruence as a central mechanism, but not as the only determinant of the overall experience.
